{{دیوان
| عنوان = بلالؓ (چمک اُٹھا جو ستارہ ترے مقدّر کا)
| ماخذ = ویکی ماخذ
| ماخذ_لنک = https://ur.wikisource.org/wiki/%D8%A8%D9%84%D8%A7%D9%84_%28%DA%86%D9%85%DA%A9_%D8%A7%D9%8F%D9%B9%DA%BE%D8%A7_%D8%AC%D9%88_%D8%B3%D8%AA%D8%A7%D8%B1%DB%81_%D8%AA%D8%B1%DB%92_%D9%85%D9%82%D8%AF%D9%91%D8%B1_%DA%A9%D8%A7%29
| مقابلہ = اقبال سائبر لائبریری (اقبال اکادمی پاکستان)
| مقابلہ_لنک = https://www.iqbalcyberlibrary.net
}}
<poem>
چمک اُٹھا جو ستارہ ترے مقدّر کا
حبَش سے تجھ کو اُٹھا کر حجاز میں لایا

ہُوئی اسی سے ترے غم کدے کی آبادی
تری غلامی کے صدقے ہزار آزادی

وہ آستاں نہ چھُٹا تجھ سے ایک دم کے لیے
کسی کے شوق میں تو نے مزے ستم کے لیے

جفا جو عشق میں ہوتی ہے وہ جفا ہی نہیں
ستم نہ ہو تو محبت میں کچھ مزا ہی نہیں

نظر تھی صورتِ سلماںؓ ادا شناس تری
شرابِ دید سے بڑھتی تھی اور پیاس تری

تجھے نظارے کا مثلِ کلیمؑ سودا تھا
اویسؓ طاقتِ دیدار کو ترستا تھا

مدینہ تیری نگاہوں کا نور تھا گویا
ترے لیے تو یہ صحرا ہی طُور تھا گویا

تری نظر کو رہی دید میں بھی حسرتِ دید
خُنَک دلے کہ تپید و دمے نیا سائید

گری وہ برق تری جانِ ناشکیبا پر
کہ خندہ زن تری ظُلمت تھی دستِ موسیٰ پر

تپش ز شعلہ گر فتند و بر دلِ تو زدند
چہ برقِ جلوہ بخاشاکِ حاصلِ تو زدند!

ادائے دید سراپا نیاز تھی تیری
کسی کو دیکھتے رہنا نماز تھی تیری

اذاں ازل سے ترے عشق کا ترانہ بنی
نماز اس کے نظارے کا اک بہانہ بنی

خوشا وہ وقت کہ یثرِب مقام تھا اس کا
خوشا وہ دور کہ دیدار عام تھا اس کا
</poem>
